\documentclass[10pt,a4paper]{amsart}
\usepackage[margin=19mm]{geometry}
\usepackage{amsmath,amssymb,mathtools}
\usepackage[scr=rsfs,cal=euler]{mathalfa}
\usepackage{xfrac}
\usepackage[unicode,hidelinks,bookmarks=false]{hyperref}
\newcommand{\ie}{i.e.\ }
\newcommand{\cf}{cf.\ }
\newcommand{\resp}{respectively\ }
\newcommand{\Subsection}{\subsection}
\providecommand{\nobreakdash}{\nobreak\hbox{-}}
\setlength{\emergencystretch}{2em}
\multlinegap0pt
\usepackage{amssymb}
\usepackage{xcolor}
\newtheorem{theorem}{Theorem}
\title{On $S$-Integral Domains and $S$-Version of Krull Intersection Theorem}
\author{Tushar Singh, Gyanendra K. Verma,, Shiv Datt Kumar}
\date{Source: arXiv:2512.20316v1 (CC BY 4.0)}
\begin{document}
\maketitle

\noindent\emph{Source and licence.} On $S$-Integral Domains and $S$-Version of Krull Intersection Theorem. Version arXiv:2512.20316v1 (CC BY 4.0), distributed by arXiv under the Creative Commons Attribution 4.0 International licence, http://creativecommons.org/licenses/by/4.0/, as recorded in the arXiv licence metadata for this exact version and shown on the abstract page https://arxiv.org/abs/2512.20316v1. Full licence notice: \texttt{licenses/arxiv-2512.20316-CC-BY-4.0.txt}.\par\smallskip

\noindent\textit{Editorial scope.} Counted unit: the article's theorem S-1r (source0.tex lines 351--357). The article's complete original proof of it is delivered with it (source0.tex lines 358--374, one \texttt{proof} environment of 17 lines). No mathematics is written, repaired or re-derived: the statement and the proof are the article's own lines, in the article's own notation, and no retained line is substituted or re-flowed.  No further source-local context is needed: every cross-reference of every retained line resolves inside the two delivered spans.  Nothing was omitted from any retained span, so the package carries no omission record.  No retained line contains a citation, so the wrapper carries no bibliography and no bibliography entry.  No retained line contains a figure environment, an image-inclusion command or drawing code, so no image is delivered and the package declares no asset dependency.  Editorial formatting only, in addition: an AMS \texttt{amsart} preamble that restates the article's own declarations and macros used by the retained lines, namely the environments \texttt{theorem} on separate counters, together with the standard packages those lines need.  The retained environments print in this document as Theorem 1 in order of appearance, whereas the article prints the same items with the numbering of its own full document.  The complete native source of the article as distributed in the arXiv e-print for this exact version is bundled unchanged as \texttt{sources/191552/source0.tex}.
	\begin{theorem}\label{S-1r}
		Let $S$ be a proper multiplicatively closed subset of $R$, and $S^{-1}R$ is a $\phi(S)$-field. Then $R$ is an $S$-field if any one of the following holds
		\begin{enumerate}
			\item $S$ is a finite set.
			\item $R$ is $S$-PID.
		\end{enumerate}
	\end{theorem}

	\begin{proof}
		\begin{enumerate}
			\item Let $J$ be an ideal of $R$. Then $S^{-1}J\subseteq S^{-1}R$ is an ideal containing the zero ideal $[0]$ of $S^{-1}R.$ Since $S^{-1}R$ is a $\phi(S)$-field, therefore, $S^{-1}J\cap \phi(S)\neq \emptyset $ or there exists $\dfrac{t}{1}\in \phi(S)$ such that $\dfrac{t}{1}(S^{-1}J)\subseteq [0].$
			\begin{description}
				\item[Case 1:] Suppose $S^{-1}J\cap \phi(S)\neq \emptyset$ and $\dfrac{a}{s}\in S^{-1}J\cap \phi(S).$ This implies that $\dfrac{a}{s}=\dfrac{j}{s_1}=\dfrac{s_2}{1}$ for some $s_1,s_2\in S$ and $j\in J.$ Consequently, there exists $u\in S$ such that $u(j-s_1s_2)=0$. Thus $j=s_1s_2$ since $S$ is proper. Thus $s_1s_2\in J\cap S\neq \emptyset$.
				\item[Case 2:] Suppose $\dfrac{t}{1}S^{-1}J\subseteq [0],$ where $\dfrac{t}{1}\in \phi(S)$. This implies that $S^{-1}J\subseteq [0]$, since $\dfrac{t}{1}$ is a unit in $S^{-1}R.$ Consequently, $\dfrac{j}{s}=\dfrac{0}{1}$ for all $s\in S$ and $j\in J.$ In particular, for $s=1,$ $\dfrac{j}{1}=\dfrac{0}{1}$, $\forall j\in J$ which implies that there exists $s_j\in S$ for all $j$ such that $s_j j=0.$ Take $u=\prod_{s\in S} s$, then $uJ\subseteq (0),$ because $u\in S$ is finite product of elements in $S.$
			\end{description}
			Thus $(0)$ is an $S$-maximal that is $R$ is an $S$-field.
			\item Let $(0)\subseteq J$ be an ideal of $R.$ Since $R$ is an $S$-PID, there exist $s\in S$ and $a\in J$ such that $sJ\subseteq \langle a\rangle.$ Thus to show $R$ is an $S$-field, it is sufficient to show that whenever $(0)\subseteq \langle a\rangle$, then $\langle a\rangle \cap S\neq \emptyset$ or there exists $t\in S$ such that $t\langle a\rangle\subseteq (0).$ 
			Let $S^{-1}R$ be $\phi(S)$-field, and $(0)\subseteq \langle a\rangle,$ where $a\in R$. Then $S^{-1}\langle a \rangle$ is an ideal of $S^{-1}R$ containing the zero ideal $[0]$ of $S^{-1}R.$ Therefore, $S^{-1}\langle a\rangle\cap\phi(S)\neq \emptyset$ or there exists $\dfrac{t}{1}\in \phi(S)$ such that $\dfrac{t}{1}S^{-1}\langle a\rangle \subseteq [0].$
			\begin{description}
				\item[Case 1:] Let $S^{-1}\langle a\rangle \cap \phi(S)\neq \emptyset$ and $\dfrac{t}{s}\in S^{-1}\langle a\rangle \cap \phi(S).$ This implies that $\dfrac{t}{s}=\dfrac{b}{s_1}=\dfrac{s_2}{1}$ for some $b\in\langle a\rangle $ and $s_1,s_2\in S.$ Consequently, $u(b-s_1s_2)=0$ for some $u\in S$. Thus $b=s_1s_2,$ since $S$ is proper. Hence $b\in \langle a\rangle \cap S\neq \emptyset.$
				\item[Case 2:] Let $\dfrac{t}{1}(S^{-1}\langle a\rangle)\subseteq [0]$. This implies that $S^{-1}\langle a\rangle \subseteq [0],$ since $\dfrac{t}{1}$ is a unit in $S^{-1}R$. Consequently, $s\langle a\rangle\subseteq (0)$ for some $s\in S$.
			\end{description}
			Thus $(0)$ is an $S$-maximal that is $R$ is $S$-field.
		\end{enumerate}
	\end{proof}

\end{document}
